\chapter{Reference: commands}
\label{ch:refman-cmd}

\begin{aside}
A \code{Cmd<Msg>} is a description of an effect to perform.
It is returned from \code{update}; the runtime performs the
effect and, if it produces a result, dispatches a message.
\end{aside}

\section{Constructors}

\begin{itemize}[leftmargin=*]
  \item \code{Cmd<Msg>\{\}} — no effect.
  \item \code{Cmd<Msg>::none()} — same; explicit form.
  \item \code{Cmd<Msg>::quit()} — exit the program.
  \item \code{Cmd<Msg>::after(d, msg)} — delayed message.
  \item \code{Cmd<Msg>::batch(a, b, ...)} — combine.
  \item \code{Cmd<Msg>::task(callable, then\_msg)} — run on a
    worker, dispatch result.
  \item \code{Cmd<Msg>::http\_get(url, then\_msg)} — async
    HTTP.
  \item \code{Cmd<Msg>::read\_file(path, then\_msg)} — async
    file read.
\end{itemize}

\section{Working example}

\begin{lstlisting}[language=C++]
#include <maya/maya.hpp>
#include <chrono>

using namespace maya;
using namespace maya::dsl;

struct CmdDemo {
    struct Model { int n = 0; std::string status; };
    struct Bump {};
    struct DelayedBump {};
    struct Quit {};
    using Msg = std::variant<Bump, DelayedBump, Quit>;

    static Model init() { return {}; }

    static auto update(Model m, Msg msg)
        -> std::pair<Model, Cmd<Msg>>
    {
        return std::visit(overload{
            [&](Bump) {
                m.n++;
                m.status = "bumped";
                return std::pair{m, Cmd<Msg>{}};
            },
            [&](DelayedBump) {
                m.status = "scheduled...";
                return std::pair{m,
                    Cmd<Msg>::after(std::chrono::seconds(1),
                                    Bump{})};
            },
            [](Quit) {
                return std::pair{Model{}, Cmd<Msg>::quit()};
            },
        }, msg);
    }

    static Element view(const Model& m) {
        return v(
            text(std::format("n = {}", m.n)) | Bold,
            text(m.status) | Dim,
            blank_,
            t<"b bump, d delayed, q quit"> | Dim
        ) | pad<1> | border_<Round>;
    }

    static auto subscribe(const Model&) -> Sub<Msg> {
        return key_map<Msg>({
            {'q', Quit{}},
            {'b', Bump{}},
            {'d', DelayedBump{}},
        });
    }
};

int main() { run<CmdDemo>({.title = "cmd"}); }
\end{lstlisting}

\section{Notes}

\begin{itemize}[leftmargin=*]
  \item Commands are values. You build them, return them,
    never invoke them yourself.
  \item Order inside \code{batch} is preserved.
  \item \code{quit} is irrevocable: returning it once tears
    the loop down.
\end{itemize}

\section{Recap}

\begin{recap}
\begin{itemize}[leftmargin=*]
  \item Effects are data; the runtime performs them.
  \item Common constructors: none, quit, after, batch, task,
    http\_get, read\_file.
\end{itemize}
\end{recap}

\section*{Workshop}

\begin{enumerate}[leftmargin=*]
  \item Implement a button that triggers two commands at
    once with \code{batch}.
  \item Use \code{task} to compute a fibonacci on a worker
    thread; show the result.
  \item Trigger a \code{quit} two seconds after a key press.
\end{enumerate}
